% !TeX root = ../papers/nc-p3.tex
\section{全有限配置での特異反対根の排除と八次判定}
\label{sec:all-finite-opposite-octics}
\begin{theorem}\label{thm:all-finite-opposite-octics}
任意の四辺有限な標準有限Frobenius窓で、反対座標が射影非同値なら
その対称二次式行列は一般階数4である。行列式は $H^0(\mathcal I_X^3(8))$ の非零元となる。
この空間が零の最初の有限辺は、非被約性や台点数にかかわらず全成分で
標準正則延長を高々一つ持ち、転置辺の三次交叉と既出の一交換対称判定でその有無を決められる。
\end{theorem}

\begin{proof}[接触スキームの四次元可換環]
一般階数1は射影同定を与え、2は既に排除した。3なら前節の
$N=UU^{\mathsf t}$、$E=X\cap V(I_2(U))$、$M(t)$ を取る。
$E$ 上の $U=pq^{\mathsf t}$ に対し、等方二次曲線上の一般点 $t_0$ での
$M(t_0)$ は可逆であり、$q^{\mathsf t}t_0$ は単元である。これにより出力空間を入力と同定し、
$M(t_0)^{-1}M(t)$ は $E$ の入力座標列に
$\lambda(t)=q^{\mathsf t}t/(q^{\mathsf t}t_0)$ 倍として作用する。
入力一次座標の $E$ 上の制限の空間 $P\subset\mathcal O_E$ は自明化で単元1を含むように取れる。
閉埋込みのアフィン座標なので $P$ は環として $\mathcal O_E$ を生成する。
$P$ は三次元空間 $W=\langle\lambda(t)\rangle$ による乗法で保たれる。
実際 $\lambda(t)=0$ なら $M(t)$ は $E$ の線型包を消すが、そのベクトル次元は3以上であり、
非零 $t$ での階数下界2に反する。$T=k[W]$ とすれば $1\in W$、$T\subset P$ である。
$\dim P=3$ なら $P=W=T=\mathcal O_E$ となり、長さ4以上の条件に反する。
よって $\dim P=4$。各 $0\ne w\in W$ の $P$ における乗法は階数2以上である。
入力座標が独立なのでこの乗法の行列は $M(t_0)^{-1}M(t)$ の転置に等しいからである。

$\dim T=3$ なら $T=W$。三つの剰余体因子があれば四次元加群上で少なくとも一つの
冪等元の階数は1である。因子が二個なら両冪等元は階数2を要し、長さ2の因子の根基元は
その二次元部分上の非零冪零作用なので階数1である。局所の三次単項代数なら三乗零の
四行行列の二乗の階数は高々1である。残りの局所型は二次元の平方零根基を持つ。
根基の全作用は共通消去空間への、その商からの写像となる。全非零元の階数が2以上なら
両空間とも二次元となり、二行二列の二次元鉛筆の全非零元が可逆になってしまう。
従って $\dim T=4$、$T=P=\mathcal O_E$、$\operatorname{length}E=4$ である。
その各局所socleは一次元である。そうでなければ超平面 $W$ がそのsocleと非自明に交わり、
正則表現で階数1の乗法が生じる。ゆえに $T$ はFrobeniusで、乗法の転置は完全跡対合により
通常の乗法と同時共役である。以下 $M(t)$ を、入出力の定数基底を取って
$T$ 上の $\lambda(t)$ 倍の行列として扱える。
\end{proof}

\begin{proof}[二つの長さ二因子への分解]
前節の二枚の長さ評価を $F$ の長さとも足し合わせると
\[
  20\le\operatorname{length}\Gamma+2\operatorname{length}E
       =\operatorname{length}\Gamma+8
       \le16+d\le20 .
\]
従って $d=4$ で、二つのSmith指数は行列式の各零点で等しい。
等方曲線を $\mathbb P^1$ と同定し、$\lambda(t)$ のその上の制限を $T$ 係数の二項二次式と読む。
これにより $E$ は二項二次式の平面 $\operatorname{Sym}^2\mathbb P^1$ の部分スキーム $\Xi$ になる。
実際三係数の比が $T$ を生成する。各閉点で二根は相異なる。$q$ の直交直線による切断で
$q^{\mathsf t}q$ が単元だからである。根の対応 $D_8\subset\Xi\times\mathbb P^1$ は
$\Xi$ の上で次数2のエタール被覆である。第二射影での構造層の表示行列は $M(t)$ である。
局所表示が $T[z]/(\lambda(z))$ だからである。よってその像は
$\mathbb P^1$ 上の長さ4の因子 $Z$ であり、$D_8$ は $Z$ 上有限自由局所階数2である。
等しい二つの正Smith指数が $a$ の点では $(k[[z]]/(z^a))^2$ を得るためである。

$Z$ の台を頂点とし、$\Xi$ の点をその異なる二根を結ぶ辺と読む。
各辺の局所スキームは対応する二つの局所因子 $Z_i\times Z_j$ に閉埋込みで入る。
二根は別々の近傍で分かれており、どちらも $Z$ を経由し、二根が元の二次式を一意に決めるからである。
各頂点での辺長さの和は $2\operatorname{length} Z_i$ である。
四頂点なら全辺は被約で、単純な次数2のグラフだから四角形となる。
三頂点なら長さは $(2,1,1)$ で、長さ2の頂点から両方の単純頂点へそれぞれ長さ2の辺があり、
残る辺はない。二頂点なら全辺が両方を結ぶので各頂点長は2となり、辺は両局所因子の全積になる。
一頂点は異なる二根に反する。従って全ての場合
\[
       \Xi=H\times H'\ \subset\operatorname{Sym}^2\mathbb P^1,
       \qquad \operatorname{length}H=\operatorname{length}H'=2,\quad H\cap H'=\varnothing .
\]
右辺の埋込みは二点の順序を忘れる射である。各場合とも全積への包含と局所長の等号を使っている。
二次式の表示は $T$ の単元倍までなので、以下その単元は一つの入力基底変更で吸収する。
\end{proof}

\begin{proof}[被約な四角形]
$H,H'$ が被約なら $T=k^4$。四つの点の二次式の係数を $q_i$ とし、入力を
$x_i$ の四座標に取れば、$N_{ij}=c_{ij}x_i x_j$、$C=(c_{ij})$ は $q_i$ の
非退化三次元対合でのGram行列である。対合は二項二次式の判別式の偏極である。
$c_{ii}\ne0$ であり、どの三つの $q_i$ も独立である。実際 $W$ の一つの非零切断が
三つの点で消えれば $T$ 上の乗法が階数1以下となる。

$E$ のある座標点 $i$ で $x_i$ を単元とする。$q_i$ と根を一つ共有する二つの隣の
$q_j$ は、$q_i^\perp$ へ射影すると異なる等方方向となる。
同じ根を持つ二次式の全鉛筆では判別式が一つの一次式の平方なので
二次元Gram行列式が零であること、また三つの点の独立性から従う。
反対の $q_k$ の射影は、同じ独立性からそのどちらの方向にも属さない。
二次元非退化平面では等方直線は自身にだけ直交する。
従って残る三点の射影は、どの二つを取っても相互積が非零である。
$N$ の $i$ 成分でのSchur補より、$D=\mathcal O_{X,i}$ で $x_jx_k=0$
（$j,k\ne i$、$j\ne k$）。従って三つの相異なる座標の積である全最大行列式はそこで消える。
$E$ の外では前節の $(I+I')J=0$ により初めから全て消える。
非零三次式の評価単射に矛盾する。
\end{proof}

\begin{proof}[二重点を含む四角形]
$H$ に二重点があるとしてそれを無限遠に置く。$H'$ は被約なら有限座標 $1,-1$、
二重点なら $0$ に置ける。射影直線の座標変更は等方二次曲線を保つのでGram形式は
基底変更と共通定数倍しか変わらない。二項座標を $(s^2,st,t^2)$ とすれば
Gram行列は、等方方向の方程式 $w_0w_2-w_1^2=0$ の行列の逆である
\[
       C=\begin{pmatrix}0&0&2\\0&-1&0\\2&0&0\end{pmatrix}.
\]
$H'$ が被約なら $T=k[\epsilon]/(\epsilon^2)\times k[\epsilon]/(\epsilon^2)$ で
入力を $(a+b\epsilon,c+d\epsilon)$ とする。二つの二次式
$(s+\epsilon t)(t-s)$、$(s+\epsilon t)(t+s)$ の三係数を入力に掛けた列は
\[
 U_1=\begin{pmatrix}-a&a&0\\-b&b-a&a\\c&c&0\\d&d+c&c\end{pmatrix}.
\]
両方が二重点なら $T=k[\epsilon,\delta]/(\epsilon^2,\delta^2)$ で入力を
$a+b\epsilon+c\delta+d\epsilon\delta$ とする。
$(s+\epsilon t)(t+\delta s)$ の三係数の列は
\[
 U_2=\begin{pmatrix}0&a&0\\0&b&a\\a&c&0\\b&d+a&c\end{pmatrix}.
\]
いずれも $N_i=U_i C U_i^{\mathsf t}$ である。
$N_i$ の二次小行列式を行・列の零始まり添字で $D_{ab,cd}$ と記し、
$A=D_{01,02}, B=D_{01,12}, F=D_{01,23}, G=D_{03,03}$ と置く。
一番目では $f=a^2c$ に対して
\[
 (af,bf,cf,df)=(-A/4,(A-B)/4,F/8,-(F+G)/8).
\]
二番目では $f=a^3$ に対して
\[
 (af,bf,cf,df)=(-A/2,-D_{01,03}/2,-D_{02,03}/2,(A+F)/2-G/4).
\]
これは表示行列の二行列式の展開である。いずれも四座標の一つが各点で単元なので
非零三次式 $f$ は $X$ 上で消え、三次評価の単射性に反する。
全場合が排除され、非零行列式の三重消滅および延長の結論は前節までの議論による。
\end{proof}
